% =====================================================================
% Clase -- Cálculo 11° -- Trimestre I -- Semana 02
% Sesiones 003 (lun 10 ago., 100 min), 004 (jue 13 ago., 55 min), 005 (vie 14 ago., 50 min)
% Tema 1 de la guía: Los sistemas numéricos
% Material del docente (no se entrega a los estudiantes).
% =====================================================================
\documentclass[11pt]{article}
\makeatletter\def\input@path{{../../../../../plantillas/estilo/}}\makeatother
\usepackage{mmcantillo}
\guiadelestudiante{../../guia-didactica/guia-periodo-I-numeros-reales}

\tipodocumento{Clase}
\asignatura{Cálculo}
\titulo{Semana 2: propiedades, comparación y representaciones}
\grado{11°}
\periodo{I}
% DBA: matematicas grado 11 · DBA 1 — Utiliza las propiedades de los números (naturales,
%      enteros, racionales y reales) y sus relaciones y operaciones para construir y comparar
%      los distintos sistemas numéricos.
% Evidencias trabajadas: (1) describe propiedades comunes y diferentes de los sistemas
%      numéricos; (3) construye representaciones de los conjuntos numéricos y establece
%      relaciones acordes con sus propiedades.
% Recurso: recursos/matematicas/Guías pedagógicas Matemáticas/08 - Modulo_Matematicas_Octavo.docx
%      (Tema 1: la propiedad distributiva como área de un rectángulo dividido en dos)

\begin{document}

\encabezado*

\begin{center}
{\renewcommand{\arraystretch}{1.3}
\begin{tabularx}{\linewidth}{@{}l l l X l@{}}
  \toprule
  \textbf{Sesión} & \textbf{Fecha} & \textbf{Min} & \textbf{Subtema} & \textbf{Entregables}\\
  \midrule
  003 & lun 10 ago., 07:50 & 100 & Propiedades de las operaciones en $\N$, $\Z$ y $\Q$ & Taller 2\\
  004 & jue 13 ago., 06:55 & 55  & Comparación de los sistemas numéricos & Quiz 1\\
  005 & vie 14 ago., 13:40 & 50  & Representaciones de los conjuntos numéricos & Tarea 2 (para el jue 20 ago.)\\
  \bottomrule
\end{tabularx}}
\end{center}

\small El lunes 17 de agosto es festivo: la Tarea 2 se entrega el jueves 20. \normalsize
Referencia: \refguia{tema:sistemas}.

% ---------------------------------------------------------------------
\seccion{Sesión 003 — lunes 10 de agosto · 07:50 · 100 min}

\textbf{Objetivo.} Identificar qué propiedades de las operaciones se cumplen en cada sistema
numérico y cuáles no, y justificarlo con contraejemplos.

\begin{center}
{\renewcommand{\arraystretch}{1.3}
\begin{tabularx}{\linewidth}{@{}l l X@{}}
  \toprule
  \textbf{Momento} & \textbf{Min} & \textbf{Actividad}\\
  \midrule
  Tarea 1 & 0–15 & Comparan respuestas en parejas; resuelve en el tablero los puntos 2d
    ($2{,}\overline{1} = \frac{19}{9}$) y 4 ($\frac{1}{7}$ y los restos).\\
  Explicación & 15–45 & Propiedades: cerradura, conmutativa, asociativa, distributiva,
    neutro, inverso aditivo e inverso multiplicativo. Construye con el grupo la tabla de
    abajo, pidiendo un contraejemplo para cada casilla que no se cumple. La distributiva se
    muestra como el área de un rectángulo de $6 \times (10 + 4)$ partido en dos.\\
  Discusión & 45–55 & ¿Por qué no se divide por cero? Si $\frac{5}{0} = b$, entonces
    $5 = 0 \cdot b = 0$, absurdo; y $\frac{0}{0}$ podría ser cualquier número.\\
  Taller 2 & 55–90 & En parejas.\\
  Cierre & 90–100 & Cada pareja dice una propiedad que se cumple en $\Q$ y no en $\Z$
    (inverso multiplicativo; cerradura de la división).\\
  \bottomrule
\end{tabularx}}
\end{center}

\textbf{Tabla que se construye con el grupo} (✓ se cumple, ✗ no se cumple):
\begin{center}
{\renewcommand{\arraystretch}{1.3}
\begin{tabular}{lccc}
  \toprule
  & $\N$ & $\Z$ & $\Q$\\
  \midrule
  Suma y multiplicación cerradas, conmutativas y asociativas & ✓ & ✓ & ✓\\
  Resta cerrada & ✗ ($3 - 5$) & ✓ & ✓\\
  División cerrada (divisor $\neq 0$) & ✗ & ✗ ($3 \div 5$) & ✓\\
  Todo número tiene opuesto (inverso aditivo) & ✗ & ✓ & ✓\\
  Todo número $\neq 0$ tiene inverso multiplicativo & ✗ & ✗ ($\frac{1}{2} \notin \Z$) & ✓\\
  \bottomrule
\end{tabular}}
\end{center}

% ---------------------------------------------------------------------
\seccion{Sesión 004 — jueves 13 de agosto · 06:55 · 55 min}

\textbf{Objetivo.} Comparar y ordenar números de distintos sistemas y reconocer la cadena
$\N \subset \Z \subset \Q$.

\begin{center}
{\renewcommand{\arraystretch}{1.3}
\begin{tabularx}{\linewidth}{@{}l l X@{}}
  \toprule
  \textbf{Momento} & \textbf{Min} & \textbf{Actividad}\\
  \midrule
  Repaso & 0–20 & Tres formas de comparar racionales: común denominador, productos
    cruzados ($\frac{a}{b} < \frac{c}{d} \iff ad < bc$ si $b, d > 0$) y expresión decimal.
    Ejemplo: $\frac{5}{7}$ y $\frac{7}{10}$: $50 > 49$, así que $\frac{5}{7} > \frac{7}{10}$.
    Con negativos: $-\frac{3}{4} > -0{,}8$ porque $-0{,}75$ está a la derecha de $-0{,}8$.\\
  Dudas & 20–30 & Preguntas sobre el Taller 2 y la tabla de propiedades.\\
  Quiz 1 & 30–50 & Individual, 20 minutos.\\
  Cierre & 50–55 & Se recogen los quices. Pregunta para el viernes: ¿«los números mayores
    que 5» son $6, 7, 8, \ldots$?\\
  \bottomrule
\end{tabularx}}
\end{center}

% ---------------------------------------------------------------------
\seccion{Sesión 005 — viernes 14 de agosto · 13:40 · 50 min}

\textbf{Objetivo.} Representar los conjuntos numéricos de varias formas (por extensión, por
comprensión, en la recta, con diagramas) y reconocer cuándo dos representaciones describen
el mismo conjunto. Es la última hora del viernes: sesión activa, en grupos.

\begin{center}
{\renewcommand{\arraystretch}{1.3}
\begin{tabularx}{\linewidth}{@{}l l X@{}}
  \toprule
  \textbf{Momento} & \textbf{Min} & \textbf{Actividad}\\
  \midrule
  Situación & 0–10 & Adaptada del ejemplo del DBA 1: en el tablero,
    (a) $\{6, 7, 8, \ldots\}$, (b) una recta con todo lo que está a la derecha de 5,
    (c) «los números mayores que cinco». Federico dice que las tres representan lo mismo.
    ¿Tiene razón? \emph{(No: (a) solo tiene naturales; (b) incluye $5{,}1$, $\frac{11}{2}$,
    etc.; (c) es ambigua hasta decir en qué conjunto trabajamos.)}\\
  Grupos & 10–35 & Grupos de 4. Cada grupo recibe un conjunto (los enteros entre $-3$ y 3;
    los racionales entre 0 y 1; los naturales pares; los enteros negativos) y lo representa
    en un pliego de cuatro maneras: extensión, comprensión, recta y diagrama de Venn
    $\N \subset \Z \subset \Q$ ubicando tres de sus elementos.\\
  Galería & 35–45 & Los carteles se pegan en la pared; cada grupo revisa el de otro y deja
    una pregunta escrita.\\
  Cierre & 45–50 & ¿Qué conjunto no se pudo escribir por extensión? (Los racionales entre 0 y
    1: son infinitos y no hay «siguiente»; se verá en el Tema 2.) Se asigna la Tarea 2 para
    el jueves 20.\\
  \bottomrule
\end{tabularx}}
\end{center}

\begin{nota}
Lleva pliegos de papel y marcadores. Si el tiempo no alcanza para la galería, cada grupo
presenta su cartel en un minuto.
\end{nota}

\end{document}
